\chapter{Phương pháp sinh câu hỏi trắc nghiệm bằng luồng đa tác tử}
\label{chap:phuong-phap}

\section{Kiến trúc hệ thống}
\label{sec:da-tac-tu}
Chất lượng của câu hỏi trắc nghiệm được sinh ra phụ thuộc trực tiếp vào cách thức thông tin được tổ chức và truyền đạt đến mô hình ngôn ngữ tại mỗi bước trong pipeline. Phần này trình bày thiết kế đối tượng trạng thái chia sẻ, luồng xử lý giữa các tác tử, vai trò và prompt của từng tác tử, cùng các kỹ thuật đảm bảo tính ổn định của hệ thống.

\begin{figure}[htbp]
    \centering
    \resizebox{\textwidth}{!}{\begin{tikzpicture}[
    >=Stealth,
    node distance=1.2cm and 1.6cm,
    % Styles
    input/.style={draw, rounded corners=3pt, fill=gray!10, minimum width=2cm, minimum height=0.9cm, font=\small, align=center},
    agent/.style={draw, rounded corners=3pt, fill=#1!12, minimum width=2.6cm, minimum height=1cm, font=\small\bfseries, align=center, line width=0.6pt},
    agent/.default={blue},
    output/.style={draw, rounded corners=3pt, fill=green!8, minimum width=2cm, minimum height=0.9cm, font=\small, align=center},
    loop/.style={draw, dashed, rounded corners=6pt, inner sep=6pt, line width=0.5pt, #1},
    loop/.default={gray!60},
    arr/.style={->, thick, color=#1},
    arr/.default={black!70},
    % Đã xóa fill=white để bỏ nền trắng cho các nhãn chữ
    lblnode/.style={font=\scriptsize, inner sep=1.5pt},
]

% ── Input ──
\node[input] (input) {Đoạn văn};

% ── Generators (3 parallel) ──
\node[agent=blue, above right=0.6cm and 2.6cm of input] (gen1) {$A_{\text{gen}}^1$\\[-2pt]\scriptsize Cấp 1};
\node[agent=blue, right=2.6cm of input] (gen2) {$A_{\text{gen}}^2$\\[-2pt]\scriptsize Cấp 2};
\node[agent=blue, below right=0.6cm and 2.6cm of input] (gen3) {$A_{\text{gen}}^3$\\[-2pt]\scriptsize Cấp 3};

% ── Classifier ──
\node[agent=orange, right=1.8cm of gen2] (cls) {$A_{\text{cls}}$\\[-2pt]\scriptsize Phân loại};

% ── Student ──
\node[agent=teal, right=1.8cm of cls] (stu) {$A_{\text{stu}}$\\[-2pt]\scriptsize Học sinh};

% ── Fixer ──
\node[agent=red, below=1cm of stu] (fix) {$A_{\text{fix}}$\\[-2pt]\scriptsize Hiệu chỉnh};

% ── Output ──
\node[output, right=1.8cm of stu] (output) {Bộ MCQ\\[-1pt]\scriptsize đạt chuẩn};

% ── Arrows: Input → Generators ──
\coordinate (bus_center) at ([xshift=1.6cm]input.east);
\coordinate (feedback_point) at ([xshift=0.8cm]input.east);

% Trục chính từ Input
\draw[thick, black!70] (input.east) -- (bus_center);
% Trục phân phối dọc
\draw[thick, black!70] (bus_center |- gen1.west) -- (bus_center |- gen3.west);
% Rẽ nhánh vào từng Generator
\draw[arr] (bus_center |- gen1.west) -- (gen1.west);
\draw[arr] (bus_center) -- (gen2.west);
\draw[arr] (bus_center |- gen3.west) -- (gen3.west);

% ── Arrows: Generators → Classifier ──
\draw[arr] (gen1.east) -| ([xshift=-0.8cm, yshift=6pt]cls.west) -- ([yshift=6pt]cls.west);
\draw[arr] (gen2.east) -- (cls.west);
\draw[arr] (gen3.east) -| ([xshift=-0.8cm, yshift=-6pt]cls.west) -- ([yshift=-6pt]cls.west);

% ── Arrow: Classifier → Student ──
\draw[arr=black!70] (cls.east) -- node[lblnode, above] {\scriptsize $\hat{l}=l$} (stu.west);

% ── Arrow: Student → Output ──
\draw[arr=black!70] (stu.east) -- node[lblnode, above] {\scriptsize pass} (output.west);

% ── Arrow: Student → Fixer ──
\draw[arr=red!60] (stu.south) -- node[lblnode, right] {\scriptsize fail} (fix.north);

% ── Arrow: Fixer → Student (loop back) ──
\draw[arr=red!60] (fix.west) -| ([xshift=-0.6cm, yshift=-4pt]stu.west) -- ([yshift=-4pt]stu.west);

% ── Loop 1: Classifier → Generators (Feedback) ──
\path (gen1.north) ++(0, 0.7cm) coordinate (loop1_top_level);

\draw[arr=orange!60] (cls.north) -- (cls.north |- loop1_top_level) -- node[lblnode, below] {\scriptsize $\hat{l}\neq l$} (feedback_point |- loop1_top_level) -- (feedback_point);

% ── Loop boxes (Background) ──
\path (feedback_point) ++(-0.2cm, 0) coordinate (loop1_left_bound);

\begin{scope}[on background layer]
    \node[loop=blue!40, fill=blue!3, fit=(loop1_left_bound)(gen1)(gen2)(gen3)(cls)(loop1_top_level), label={[font=\scriptsize\itshape, blue!60, yshift=2pt]above:Vòng lặp 1: Kiểm soát cấp độ nhận thức}] {};
    \node[loop=red!40, fill=red!3, fit=(stu)(fix), label={[font=\scriptsize\itshape, red!60]below:Vòng lặp 2: Kiểm định tính giải được}] {};
\end{scope}

\end{tikzpicture}}
    \caption{Kiến trúc hệ thống đa tác tử với hai vòng lặp kiểm soát chất lượng: vòng lặp phân loại cấp độ nhận thức và vòng lặp kiểm định tính giải được của câu hỏi}
    \label{fig:multiagent-architecture}
\end{figure}

Toàn bộ pipeline suy luận được điều phối thông qua một đối tượng trạng thái chia sẻ \texttt{WorkflowState}, sử dụng khung điều phối \textbf{LangGraph} \cite{langgraph2024} để xây dựng đồ thị có hướng. Mỗi trạng thái $s$ gồm hai nhóm trường: trường đầu vào $\mathcal{I}$ được khởi tạo từ dữ liệu bài đọc, và trường trung gian $\mathcal{O}$ được cập nhật bởi các tác tử trong quá trình xử lý:

\begin{equation}
\mathcal{I} = \{\textit{content}, \textit{n}, \textit{max\_loops}\}
\end{equation}
\begin{equation}
\mathcal{O} = \{\textit{questions}\}
\end{equation}

Pipeline hoạt động theo hai vòng lặp kiểm soát chất lượng (Hình~\ref{fig:multiagent-architecture}). Vòng lặp thứ nhất đảm bảo câu hỏi đúng cấp độ nhận thức: Tác tử \textit{Sinh câu hỏi} ($A_{\text{gen}}^l$) tạo câu hỏi cho cấp độ $l$, sau đó Tác tử \textit{Phân loại} ($A_{\text{cls}}$) xác minh cấp độ; câu hỏi bị từ chối sẽ được sinh lại ở vòng tiếp theo. Vòng lặp thứ hai đảm bảo câu hỏi có thể giải được: Tác tử \textit{Học sinh} ($A_{\text{stu}}$) thử giải câu hỏi, nếu thất bại thì Tác tử \textit{Hiệu chỉnh} ($A_{\text{fix}}$) sửa phương án rồi đưa lại cho Tác tử \textit{Học sinh} kiểm định. Mỗi câu hỏi có số lần thử lại tối đa có thể cấu hình; nếu vượt quá giới hạn, câu hỏi sẽ được đánh dấu lỗi và loại khỏi kết quả cuối cùng.

\section{Luồng xử lý đa tác tử}

\subsection{Tác tử \textit{Sinh câu hỏi}}
Hệ thống triển khai ba tác tử \textit{sinh câu hỏi} hoạt động song song, mỗi tác tử chịu trách nhiệm cho một cấp độ nhận thức $l \in \{1, 2, 3\}$ theo thang đo Bloom. Mỗi tác tử nhận đoạn văn bản nguồn, số lượng câu hỏi cần sinh, cùng hai tập mẫu tham chiếu từ các vòng lặp trước:

\begin{equation}
A_{\text{gen}}^l: (\textit{content}, n_l, \mathcal{P}_l, \mathcal{N}_l) \longrightarrow \{q_1, q_2, \ldots, q_{n_l}\}
\end{equation}

\noindent trong đó $n_l$ là số câu hỏi cần sinh cho cấp độ $l$, $\mathcal{P}_l$ là tập mẫu tích cực và $\mathcal{N}_l$ là tập mẫu tiêu cực. Mỗi câu hỏi $q_i$ được biểu diễn dưới dạng bộ:

\begin{equation}
q_i = (\textit{id}, \textit{question}, \textit{options} \in \mathbb{S}^4, \textit{correct\_idx} \in \{0,1,2,3\}, \textit{level} = l, \textit{status})
\end{equation}

Prompt của mỗi tác tử được thiết kế chuyên biệt cho cấp độ nhận thức tương ứng: tác tử \textit{Cấp 1} yêu cầu câu hỏi trích xuất thông tin trực tiếp từ văn bản; tác tử \textit{Cấp 2} đòi hỏi câu hỏi suy luận ngầm, đáp án phải được diễn đạt lại hoàn toàn so với văn bản gốc; tác tử \textit{Cấp 3} yêu cầu câu hỏi lập luận đa bước theo chuẩn LSAT/GMAT, kết nối các thông tin rải rác trong nhiều đoạn văn.

Một kỹ thuật quan trọng trong prompt của tác tử \textit{sinh câu hỏi} là cơ chế học từ mẫu tham chiếu. Trước khi sinh câu hỏi, mô hình nhận được danh sách các mẫu tích cực $\mathcal{P}_l$ (câu hỏi đã đạt chuẩn từ các vòng trước) và các mẫu tiêu cực $\mathcal{N}_l$ (câu hỏi đã bị từ chối vì sai cấp độ). Prompt hướng dẫn mô hình mô phỏng phong cách và độ khó của mẫu tích cực, đồng thời tránh lặp lại các lỗi từ mẫu tiêu cực. Cơ chế này hoạt động như một dạng few-shot learning có phản hồi, trong đó chất lượng câu hỏi được cải thiện dần qua từng vòng lặp. Hình~\ref{fig:agent-generator} minh họa luồng xử lý bên trong tác tử.

\begin{figure}[htbp]
    \centering
    \resizebox{0.85\textwidth}{!}{% Agent: Generator (Sinh câu hỏi)
\begin{tikzpicture}[
    >=Stealth,
    node distance=0.8cm,
    box/.style={draw, rounded corners=2pt, minimum width=2.4cm, minimum height=0.7cm, font=\small, align=center, line width=0.5pt},
    io/.style={box, fill=gray!8},
    proc/.style={box, fill=blue!10},
    lbl/.style={font=\scriptsize, text=gray!70},
]

% ── Inputs ──
\node[io] (content) {$content$};
\node[io, below=0.35cm of content] (n) {$n_l$};
\node[io, below=0.35cm of n] (pos) {$\mathcal{P}_l$\; \scriptsize mẫu tích cực};
\node[io, below=0.35cm of pos] (neg) {$\mathcal{N}_l$\; \scriptsize mẫu tiêu cực};

% ── Input group label ──
\node[lbl, above=0.1cm of content] {Đầu vào};

% ── LLM Processing ──
\node[proc, right=2cm of n, yshift=-0.3cm, minimum height=1.8cm, minimum width=3cm, fill=blue!12] (llm) {\textbf{sLLM}\\[2pt]\scriptsize Prompt chuyên biệt\\[-1pt]\scriptsize cho cấp độ $l$\\[2pt]\scriptsize + CoT reasoning};

% ── Output ──
\node[io, right=2cm of llm, minimum height=1.2cm] (out) {$\{q_1, \ldots, q_{n_l}\}$\\[1pt]\scriptsize id, question,\\[-1pt]\scriptsize options, answer};

% ── Output label ──
\node[lbl, above=0.1cm of out] {Đầu ra};

% ── Arrows ──
\draw[->, thick, gray!70] (content.east) -- ++(0.3,0) |- ([yshift=4pt]llm.west);
\draw[->, thick, gray!70] (n.east) -- (n.east -| llm.west);
\draw[->, thick, gray!70] (pos.east) -- ++(0.3,0) |- ([yshift=-4pt]llm.west);
\draw[->, thick, gray!70] (neg.east) -- ++(0.6,0) |- ([yshift=-10pt]llm.west);
\draw[->, thick, blue!60] (llm.east) -- (out.west);

% ── Surrounding box ──
\begin{scope}[on background layer]
    \node[draw, dashed, rounded corners=5pt, inner sep=10pt, fit=(content)(n)(pos)(neg)(llm)(out), blue!30, fill=blue!2, label={[font=\small\bfseries, blue!70]above:Tác tử Sinh câu hỏi $(A_{\text{gen}}^l)$}] {};
\end{scope}

\end{tikzpicture}
}
    \caption{Luồng xử lý bên trong Tác tử Sinh câu hỏi: nhận văn bản nguồn và mẫu tham chiếu, sinh câu hỏi qua sLLM với prompt chuyên biệt theo cấp độ}
    \label{fig:agent-generator}
\end{figure}

\subsection{Tác tử \textit{Phân loại}}
Tác tử \textit{Phân loại} có nhiệm vụ xác minh rằng mỗi câu hỏi được sinh ra phù hợp với cấp độ nhận thức mong muốn:

\begin{equation}
A_{\text{cls}}: (\textit{content}, \{q_1, \ldots, q_m\}) \longrightarrow \{(\textit{id}_i, \hat{l}_i)\}_{i=1}^{m}
\end{equation}

\noindent trong đó $\hat{l}_i$ là cấp độ nhận thức do tác tử dự đoán cho câu hỏi $q_i$.

Để chống thiên lệch vị trí, trước khi phân loại, các câu hỏi được xáo trộn ngẫu nhiên và gán các mã định danh tạm thời $\{1, 2, \ldots, m\}$, khác với mã định danh toàn cục của câu hỏi. Cơ chế này ngăn mô hình xác nhận cấp độ dựa trên thứ tự xuất hiện hoặc mã định danh thay vì phân tích nội dung thực tế.

Kết quả phân loại trực tiếp cập nhật hai tập mẫu tham chiếu cho tác tử \textit{sinh câu hỏi} ở vòng lặp tiếp theo. Với mỗi câu hỏi $q$ ở cấp độ mục tiêu $l$:

\begin{equation}
\mathcal{P}_l^{(t+1)} = \mathcal{P}_l^{(t)} \cup \{q.\textit{question} \mid \hat{l}(q) = l\}
\end{equation}
\begin{equation}
\mathcal{N}_l^{(t+1)} = \mathcal{N}_l^{(t)} \cup \{q.\textit{question} \mid \hat{l}(q) \neq l\}
\end{equation}

\noindent trong đó $\mathcal{P}_l^{(t)}$ và $\mathcal{N}_l^{(t)}$ lần lượt là tập mẫu tích cực và tiêu cực tại vòng lặp $t$. Câu hỏi đạt chuẩn ($\hat{l} = l$) được chuyển sang Tác tử \textit{Học sinh}; câu hỏi bị từ chối ($\hat{l} \neq l$) được đánh dấu thất bại và nút điều phối sẽ yêu cầu tác tử \textit{sinh câu hỏi} tạo lại số lượng thiếu hụt. Hình~\ref{fig:agent-classifier} minh họa luồng xử lý bên trong tác tử.

\begin{figure}[htbp]
    \centering
    \resizebox{0.9\textwidth}{!}{% Agent: Classifier (Phân loại)
\begin{tikzpicture}[
    >=Stealth,
    node distance=0.8cm,
    box/.style={draw, rounded corners=2pt, minimum width=2.2cm, minimum height=0.7cm, font=\small, align=center, line width=0.5pt},
    io/.style={box, fill=gray!8},
    proc/.style={box, fill=orange!10},
    decision/.style={diamond, draw, aspect=2, inner sep=1.5pt, font=\small, fill=orange!8, line width=0.5pt},
    lbl/.style={font=\scriptsize, text=gray!70},
]

% ── Input ──
\node[io] (input) {$\{q_1, \ldots, q_m\}$\\[1pt]\scriptsize + $content$};
\node[lbl, above=0.1cm of input] {Đầu vào};

% ── Shuffle ──
\node[proc, right=1.4cm of input] (shuffle) {\scriptsize Xáo trộn\\[-1pt]\scriptsize + gán ID tạm};

% ── LLM ──
\node[proc, right=1.2cm of shuffle, minimum height=1.2cm, minimum width=2.6cm, fill=orange!12] (llm) {\textbf{sLLM}\\[2pt]\scriptsize Phân loại Bloom\\[-1pt]\scriptsize + CoT reasoning};

% ── Decision ──
\node[decision, right=1.2cm of llm] (dec) {\scriptsize $\hat{l}=l$\,?};

% ── Outputs ──
\node[io, above right=0.6cm and 1.2cm of dec, fill=green!8] (pass) {$\mathcal{P}_l$\\[1pt]\scriptsize Đạt chuẩn};
\node[io, below right=0.6cm and 1.2cm of dec, fill=red!8] (fail) {$\mathcal{N}_l$\\[1pt]\scriptsize Sinh lại};

% ── Arrows ──
\draw[->, thick, gray!70] (input.east) -- (shuffle.west);
\draw[->, thick, gray!70] (shuffle.east) -- (llm.west);
\draw[->, thick, orange!60] (llm.east) -- (dec.west);
\draw[->, thick, green!50!black] (dec.north east) -- node[lbl, above left, font=\scriptsize] {Có} (pass.south west);
\draw[->, thick, red!60] (dec.south east) -- node[lbl, below left, font=\scriptsize] {Không} (fail.north west);

% ── Surrounding box ──
\begin{scope}[on background layer]
    \node[draw, dashed, rounded corners=5pt, inner sep=10pt, fit=(input)(shuffle)(llm)(dec)(pass)(fail), orange!30, fill=orange!2, label={[font=\small\bfseries, orange!70]above:Tác tử Phân loại $(A_{\text{cls}})$}] {};
\end{scope}

\end{tikzpicture}
}
    \caption{Luồng xử lý bên trong Tác tử Phân loại: xáo trộn câu hỏi chống thiên lệch, phân loại cấp độ Bloom, và tách kết quả thành tập đạt chuẩn và tập sinh lại}
    \label{fig:agent-classifier}
\end{figure}

\subsection{Tác tử \textit{Học sinh}}
Tác tử \textit{Học sinh} mô phỏng quá trình một học sinh giải từng câu hỏi, giúp phát hiện các câu hỏi mơ hồ, có nhiều đáp án đúng, hoặc có cấu trúc kém:

\begin{equation}
A_{\text{stu}}: (\textit{content}, \{q_1, \ldots, q_m\}) \longrightarrow \{(\textit{id}_i, C_i, \textit{reason}_i)\}_{i=1}^{m}
\end{equation}

\noindent trong đó $C_i \subseteq \{A, B, C, D\} \cup \{\texttt{NONE}\}$ là tập hợp các phương án được tác tử lựa chọn cho câu hỏi $q_i$.

Một quyết định thiết kế quan trọng là bài kiểm tra được đưa cho tác tử dưới dạng trắc nghiệm nhiều lựa chọn đúng (Select All That Apply) thay vì trắc nghiệm một đáp án. Trong định dạng trắc nghiệm truyền thống chỉ cho phép chọn một đáp án, mô hình ngôn ngữ có xu hướng chọn phương án ``trông hợp lý nhất'' mà không đánh giá kỹ từng phương án một cách độc lập. Ngược lại, định dạng nhiều lựa chọn đúng buộc mô hình phải đánh giá từng phương án dựa trên đoạn văn nguồn, từ đó phát hiện chính xác hơn các trường hợp: câu hỏi có nhiều phương án đều đúng (tính đơn nghĩa bị vi phạm), không phương án nào đúng (đáp án sai), hoặc phương án nhiễu quá dễ loại trừ.

Câu hỏi được xác nhận đạt chuẩn khi và chỉ khi tác tử lựa chọn đúng một phương án duy nhất trùng với đáp án chính xác:

\begin{equation}
\textit{passed}(q) = \begin{cases}
\textit{true} & \text{nếu } C(q) = \{c^*\} \\
\textit{false} & \text{ngược lại}
\end{cases}
\end{equation}

\noindent trong đó $c^*$ là ký hiệu phương án đúng của câu hỏi $q$. Nếu tác tử trả lời sai, câu hỏi được chuyển tới Tác tử \textit{Hiệu chỉnh} cùng với lý do giải thích. Hình~\ref{fig:agent-student} minh họa luồng xử lý bên trong tác tử.

\begin{figure}[htbp]
    \centering
    \resizebox{0.85\textwidth}{!}{% Agent: Student (Học sinh)
\begin{tikzpicture}[
    >=Stealth,
    node distance=0.8cm,
    box/.style={draw, rounded corners=2pt, minimum width=2.2cm, minimum height=0.7cm, font=\small, align=center, line width=0.5pt},
    io/.style={box, fill=gray!8},
    proc/.style={box, fill=teal!10},
    decision/.style={diamond, draw, aspect=2, inner sep=1.5pt, font=\small, fill=teal!8, line width=0.5pt},
    lbl/.style={font=\scriptsize, text=gray!70},
]

% ── Input ──
\node[io] (input) {$\{q_1, \ldots, q_m\}$\\[1pt]\scriptsize + $content$};
\node[lbl, above=0.1cm of input] {Đầu vào};

% ── LLM ──
\node[proc, right=1.5cm of input, minimum height=1.4cm, minimum width=3cm, fill=teal!12] (llm) {\textbf{sLLM}\\[2pt]\scriptsize Giải MCQ dạng\\[-1pt]\scriptsize ``Select All That Apply''\\[-1pt]\scriptsize + CoT reasoning};

% ── Decision ──
\node[decision, right=1.3cm of llm] (dec) {\scriptsize $C=\{c^*\}$\,?};

% ── Outputs ──
\node[io, above right=0.6cm and 1cm of dec, fill=green!8] (pass) {\scriptsize Đạt chuẩn\\[-1pt]\scriptsize $\to$ Kết quả};
\node[io, below right=0.6cm and 1cm of dec, fill=red!8] (fail) {\scriptsize$(C_i, reason_i)$\\[-1pt]\scriptsize $\to A_{\text{fix}}$};

% ── Arrows ──
\draw[->, thick, gray!70] (input.east) -- (llm.west);
\draw[->, thick, teal!60] (llm.east) -- (dec.west);
\draw[->, thick, green!50!black] (dec.north east) -- node[above left, font=\scriptsize] {Đúng} (pass.south west);
\draw[->, thick, red!60] (dec.south east) -- node[below left, font=\scriptsize] {Sai} (fail.north west);

% ── Surrounding box ──
\begin{scope}[on background layer]
    \node[draw, dashed, rounded corners=5pt, inner sep=10pt, fit=(input)(llm)(dec)(pass)(fail), teal!30, fill=teal!2, label={[font=\small\bfseries, teal!70]above:Tác tử Học sinh $(A_{\text{stu}})$}] {};
\end{scope}

\end{tikzpicture}
}
    \caption{Luồng xử lý bên trong Tác tử Học sinh: giải MCQ dạng nhiều lựa chọn đúng, phân luồng câu hỏi đạt chuẩn và câu hỏi cần hiệu chỉnh kèm lý do}
    \label{fig:agent-student}
\end{figure}

\subsection{Tác tử \textit{Hiệu chỉnh}}
Tác tử \textit{Hiệu chỉnh} tiếp nhận các câu hỏi mà Tác tử \textit{Học sinh} đã trả lời sai, cùng với lý do giải thích:

\begin{equation}
A_{\text{fix}}: (\textit{content}, \{(q_i, C_i, \textit{reason}_i)\}) \longrightarrow \{(\textit{id}_i, \textit{options}_i', \textit{correct\_idx}_i')\}
\end{equation}

Ràng buộc thiết kế cốt lõi của tác tử này là chỉ được phép sửa đổi các phương án trả lời và đáp án đúng, trong khi bắt buộc giữ nguyên thân câu hỏi và cấp độ nhận thức. Ràng buộc này ngăn chặn hệ thống tạo ra câu hỏi hoàn toàn mới thay vì sửa chữa câu hỏi hiện có. Tác tử sử dụng lý do giải thích từ Tác tử \textit{Học sinh} để xác định nguyên nhân lỗi, chẳng hạn phương án nhiễu quá hiển nhiên, nhiều phương án cùng đúng, hoặc đáp án chính thức không chính xác, và thực hiện sửa chữa có mục tiêu. Sau khi sửa chữa, câu hỏi được đưa trở lại Tác tử \textit{Học sinh} để tái kiểm định. Hình~\ref{fig:agent-fixer} minh họa luồng xử lý bên trong tác tử.

\begin{figure}[htbp]
    \centering
    \resizebox{0.85\textwidth}{!}{% Agent: Fixer (Hiệu chỉnh)
\begin{tikzpicture}[
    >=Stealth,
    node distance=0.8cm,
    box/.style={draw, rounded corners=2pt, minimum width=2.2cm, minimum height=0.7cm, font=\small, align=center, line width=0.5pt},
    io/.style={box, fill=gray!8},
    proc/.style={box, fill=red!10},
    lbl/.style={font=\scriptsize, text=gray!70},
    constraint/.style={draw, rounded corners=2pt, fill=yellow!8, font=\scriptsize, align=center, minimum width=2.4cm, line width=0.4pt},
]

% ── Inputs ──
\node[io] (q) {$q_i$\; \scriptsize câu hỏi lỗi};
\node[io, below=0.35cm of q] (reason) {$reason_i$\\[-1pt]\scriptsize Lý do từ $A_{\text{stu}}$};
\node[io, below=0.35cm of reason] (choices) {$C_i$\\[-1pt]\scriptsize Đáp án đã chọn};
\node[lbl, above=0.1cm of q] {Đầu vào};

% ── LLM ──
\node[proc, right=1.8cm of reason, minimum height=1.4cm, minimum width=3cm, fill=red!12] (llm) {\textbf{sLLM}\\[2pt]\scriptsize Phân tích lỗi\\[-1pt]\scriptsize + Sửa phương án\\[-1pt]\scriptsize + CoT reasoning};

% ── Constraint box ──
\node[constraint, above=0.5cm of llm] (cst) {\textbf{Ràng buộc:}\\[-1pt] Giữ nguyên thân câu hỏi\\[-1pt] Chỉ sửa options \& answer};

% ── Output ──
\node[io, right=1.8cm of llm, minimum height=1cm] (out) {$options'_i$\\[1pt]$correct\_idx'_i$\\[-1pt]\scriptsize $\to A_{\text{stu}}$};
\node[lbl, above=0.1cm of out] {Đầu ra};

% ── Arrows ──
\draw[->, thick, gray!70] (q.east) -- ++(0.3,0) |- ([yshift=4pt]llm.west);
\draw[->, thick, gray!70] (reason.east) -- (llm.west);
\draw[->, thick, gray!70] (choices.east) -- ++(0.3,0) |- ([yshift=-4pt]llm.west);
\draw[->, thick, red!60] (llm.east) -- (out.west);
\draw[->, dashed, gray!50] (cst.south) -- (llm.north);

% ── Surrounding box ──
\begin{scope}[on background layer]
    \node[draw, dashed, rounded corners=5pt, inner sep=10pt, fit=(q)(reason)(choices)(llm)(out)(cst), red!30, fill=red!2, label={[font=\small\bfseries, red!70]above:Tác tử Hiệu chỉnh $(A_{\text{fix}})$}] {};
\end{scope}

\end{tikzpicture}
}
    \caption{Luồng xử lý bên trong Tác tử Hiệu chỉnh: nhận câu hỏi lỗi kèm lý do từ Tác tử Học sinh, chỉ sửa phương án và đáp án, đưa lại để tái kiểm định}
    \label{fig:agent-fixer}
\end{figure}

\section{Một số kỹ thuật mở rộng}

\subsection{Suy luận chuỗi tư duy}
Toàn bộ các tác tử trong pipeline đều áp dụng kỹ thuật suy luận chuỗi tư duy (CoT - Chain-of-Thought) trước khi đưa ra kết quả. Cụ thể, mỗi prompt đều yêu cầu mô hình trình bày chuỗi lập luận thông qua trường \texttt{Reason} trong định dạng đầu ra bắt buộc, trước khi đưa ra câu trả lời hoặc quyết định cuối cùng.

Thiết kế này được thể hiện nhất quán trong prompt của từng tác tử:

\begin{itemize}
    \item \textbf{Tác tử \textit{Sinh câu hỏi}:} Trường \texttt{Reason} yêu cầu mô hình giải thích chiến lược tạo câu hỏi và các phương án trước khi sinh nội dung.
    \item \textbf{Tác tử \textit{Phân loại}:} Trường \texttt{Reason} yêu cầu mô hình phân tích cấp độ nhận thức dựa trên bằng chứng cụ thể trước khi đưa ra phán quyết.
    \item \textbf{Tác tử \textit{Học sinh}:} Trường \texttt{Reason} yêu cầu mô hình trình bày chứng minh từng bước trong ba câu, giải thích lý do chọn hoặc loại từng phương án.
    \item \textbf{Tác tử \textit{Hiệu chỉnh}:} Trường \texttt{Reason} yêu cầu mô hình phân tích nguyên nhân lỗi của phương án gốc trước khi đề xuất bản sửa.
\end{itemize}

Kỹ thuật suy luận chuỗi tư duy đặc biệt quan trọng khi sử dụng các sLLM, vì các mô hình có quy mô nhỏ thường gặp khó khăn trong các tác vụ đòi hỏi suy luận nhiều bước. Bằng cách buộc mô hình trình bày rõ quá trình tư duy trước khi đưa ra kết quả, kỹ thuật này giúp giảm lỗi logic, tăng tính nhất quán của đầu ra, đồng thời cho phép truy vết và gỡ lỗi khi cần thiết.

\subsection{Định dạng đầu ra của tác tử}
Toàn bộ pipeline sử dụng đầu ra dạng văn bản thuần có cấu trúc thay vì yêu cầu mô hình trả về định dạng JSON. Cụ thể, mỗi tác tử trả về kết quả theo cấu trúc khóa-giá trị cố định như được minh họa trong Hình~\ref{fig:agent_output_format}:

\begin{figure}[htbp]
    \centering
    \begin{minipage}{0.5\textwidth}
\begin{lstlisting}[basicstyle=\ttfamily\small, backgroundcolor=\color{white}, numbers=none, frame=tb, belowskip=0pt]
ID: [ID]
Reason: [...]
Question: [text]
A: [option A]
B: [option B]
C: [option C]
D: [option D]
Answer: [A|B|C|D]
\end{lstlisting}
    \end{minipage}
    \vspace{-0.3cm}
    \caption{Cấu trúc định dạng đầu ra chuẩn của tác tử \textit{Sinh câu hỏi}}
    \label{fig:agent_output_format}
\end{figure}

Quyết định thiết kế này xuất phát từ thực tế rằng các sLLM tỏ ra kém hiệu quả và không ổn định khi phải tuân thủ định dạng JSON nghiêm ngặt. Các lỗi thường gặp bao gồm: thiếu hoặc thừa dấu ngoặc, thoát ký tự đặc biệt sai (đặc biệt với dấu ngoặc kép trong nội dung câu hỏi), cắt ngắn đầu ra giữa chừng khi chuỗi JSON quá dài, và sinh ra JSON không hợp lệ về mặt cú pháp. Những lỗi này xảy ra với tần suất đáng kể ở các sLLM có quy mô dưới 15 tỷ tham số, khiến toàn bộ đầu ra của một lần gọi mô hình bị mất khi quá trình phân tích cú pháp JSON thất bại.

Ngược lại, định dạng văn bản thuần mang lại hai lợi thế chính. Thứ nhất, mô hình ngôn ngữ sinh văn bản theo cơ chế tự hồi quy, tức sinh từng token một theo trình tự, nên cấu trúc khóa-giá trị tuần tự phù hợp tự nhiên hơn với cơ chế hoạt động của mô hình so với cấu trúc lồng nhau của JSON. Thứ hai, thiết kế này giúp pipeline chịu lỗi tốt hơn: ngay cả khi mô hình thêm các ký hiệu định dạng ngoài ý muốn như tiêu đề markdown, dấu đầu dòng hay ký hiệu in đậm, hệ thống phân tích cú pháp vẫn có thể trích xuất được dữ liệu từ đầu ra.